
\newcommand{\nomedoc}{StudioDiFattibilita}
\newcommand{\versione}{0.4}
\newcommand{\versioneglossario}{1.0}
\newcommand{\versionenormeprogetto}{1.0}
\newcommand{\nomefile}{StudioDiFattibilita-\versione.pdf}
\newcommand{\datacreazione}{7 Dicembre 2010}
\newcommand{\datamodifica}{10 Dicembre 2010}
\newcommand{\stato}{formale}
\newcommand{\uso}{interno}
\newcommand{\redazione}{Lovato Daniele, Palazzin Alberto}
\newcommand{\verifica}{}
\newcommand{\approvazione}{Valter}
\newcommand{\distribuzione}{
VT.G \\
& Prof. Vardanega Tullio\\
& Prof. Cardin Riccardo }

% FUNZIONI TIPOGRAFICHE
\newcommand{\co}{\texttt} % courier
\newcommand{\bo}{\textbf} % bold
\newcommand{\pr}{\par\medskip} % paragrafo spaziato
\newcommand{\sca}{\textsc} % small caps

\documentclass[a4paper,12pt]{report}
% 10pt,11pt,12pt
% titlepage, notitlepage -> per dare inizio o no ad una nuova pagina dopo titolo
% twoside -> per dire se fronte-retro
\usepackage[latin1]{inputenc}
% per caratteri accentati
\usepackage[italian]{babel}
% per regole sintattiche italiane
\usepackage[bookmarks=true, pdfborder={0 0 0 0}]{hyperref}
% per collegamenti ipertestuali
\usepackage{graphicx}
% per inserimento immagini

% \usepackage{enumerate}
% per personalizzare elenchi puntati

\usepackage[hmargin=2cm]{geometry} %margine 2 cm
%\geometry{options varie}

% comandi per gestire meglio header e footer
\usepackage{fancyhdr}  % header e footer
\usepackage{totpages}
\pagestyle{fancy}
\renewcommand{\headrulewidth}{0.4pt}
\renewcommand{\footrulewidth}{0.4pt}

\setlength{\headheight}{1.2cm} % NON TOCCARE
\setlength{\voffset}{-1.5cm} % NON TOCCARE
\setlength{\textheight}{666pt} % NON TOCCARE
\setlength{\footskip}{60pt}
\setlength{\parindent}{0pt} % INDENTAZIONE

\lhead{\nomedoc\  (ver. \versione)}
\chead{}
\rhead{\includegraphics[height=1cm]{img/netmus.png}}
\lfoot{\includegraphics[height=0.8cm]{img/logo.png}}
\cfoot{}
\rfoot{\thepage}

\usepackage{titlesec}
\titleformat{\chapter}{\normalfont\huge\bfseries}
{\thechapter}{20pt}{\Huge}

\usepackage{rotating}   % PER TABELLE E AMBIENTI RUOTATI
\usepackage{array}
\usepackage{color}
\usepackage{colortbl}  % VARIE PER GESTIRE I COLORI
\definecolor{Orange}{RGB}{255,127,0}   % ARANCIO ACCES0
\definecolor{orange}{RGB}{255,207,80}  % ARANCIO TENUE

\addtocontents{toc}{\protect\thispagestyle{fancy}}  % PER INDICI CON + PAGINE
\usepackage[font=it]{caption}    % PER RENDERE CORSIVE LE DIDASCALIE
\usepackage{eurosym}  % PER SIMBOLO EURO

% \usepackage{listings}   per codice sorgente

\author{VT.G - Valter Texas Group}

\begin{document}

\pagenumbering{Roman} % INIZIO NUMERAZIONE ARABA

\vspace*{1cm}
\begin{center}

\begin{LARGE} \sca{Federico Baron} \end{LARGE}\\
\vspace{0.5cm}
\begin{Large}
\emph{fede.baron.89@gmail.com} \end{Large}\\
\vspace*{1cm} \includegraphics[width=5cm]{img/logo.png}\\
\vspace{0.5cm}
\begin{Large} \emph{``Comunicazione Aumentata/Alternativa per Giovani Ospiti
della Terapia Intensiva Pediatrica''} \end{Large}\\
\vspace{3cm}
\begin{Large} \sca{\nomedoc} \end{Large}\\
\end{center}
\vspace{1cm}

% INFORMAZIONI DOCUMENTO
\begin{center}
\begin{tabular}{r|l}
\hline & \\
\bo{Nome} & \nomefile \\
\bo{Versione attuale} & \versione \\
\bo{Data creazione} & \datacreazione \\
\bo{Data ultima modifica} & \datamodifica \\
\bo{Redazione} & \redazione \\
& \\\hline
\end{tabular}
\end{center}
\newpage

% REGISTRO MODIFICHE
\section*{Registro delle modifiche}
\begin{tabular}{lll}
% REVISIONI DEL DOCUMENTO
% DALLA PIU' NUOVA ALLA PIU' VECCHIA

% TEMPLATE PER REVISIONE

\bo{Data:} 10/12/2010 &
\bo{Versione:} 0.4 &
\bo{Autore:} Palazzin Alberto\\
\hline\\
\multicolumn{3}{p{470px}}{ Correzione errori grammaticali.}\\
\\

\bo{Data:} 09/12/2010 &
\bo{Versione:} 0.3 &
\bo{Autore:} Lovato Daniele\\
\hline\\
\multicolumn{3}{p{470px}}{ Aggiunta Sezione "Componenti del Progetto e Dettagli Tecnologie".}\\ \\

\bo{Data:} 09/12/2010 &
\bo{Versione:} 0.2 &
\bo{Autore:} Palazzin Alberto\\
\hline\\
\multicolumn{3}{p{470px}}{ Aggiunta delle sottolineature per il Glossario.}\\ \\

\bo{Data:} 07/12/2010 &
\bo{Versione:} 0.1 &
\bo{Autore:} Palazzin Alberto\\
\hline\\
\multicolumn{3}{p{470px}}{ Stesura prima versione di StudioDiFattibilita.}\\ \\

\end{tabular}

% INDICE
\tableofcontents
\thispagestyle{fancy} % per lo stile di header e footer


\chapter*{Sommario}
\thispagestyle{fancy} % serve perche' nelle pagine di inizio Chapter esca headere footer
Dopo aver letto, discusso e analizzato i vari capitolati d'appalto, \`e
stato stilato il documento dello studio di fattibilit\`a per riassumere e
documentare le scelte effettuate. In sintesi \`e stato selezionato il capitolato
C2 NetMus come futuro progetto da implementare, le motivazioni sono riportate dal
Capitolo 2 in poi.

\pagenumbering{arabic} % INIZIO NUMERAZIONE NORMALE
\rfoot{\thepage\ di \pageref{TotPages}}
\addcontentsline{toc}{chapter}{Sommario}

\chapter{Introduzione}
\thispagestyle{fancy} % serve perche' nelle pagine di inizio Chapter esca header e footer

\section{Scopo del documento}
Il presente documento ha lo scopo di valutare la fattibilit\`a di progettazione
e sviluppo di ogni capitolato d'appalto rispetto alle tecnologie e alle tempistiche richieste. Nello specifico vengono descritte le motivazioni per le quali il gruppo ha scelto il capitolato C2 piuttosto che altri.


\section{Scopo del prodotto}
Il progetto \underline{NetMus} nasce con lo scopo di realizzare un sistema
software basato su \underline{cloud} \underline{computing}, per memorizzare
informazioni di brani musicali in profili utente online.\\ Tali informazioni vengono estratte da
dispositivi musicali o di archiviazione \underline{USB} al momento della loro connessione.

\section{Glossario}
Il Glossario \`e definito con un documento a parte
(\emph{Glossario-\versioneglossario.pdf}). Tutti i termini caratterizzati da
\underline{questa sottolineatura} sono ivi definiti.\\
Verr\`a sottolineata solamente la prima occorrenza di ciascun
termine presente nel Glossario, per non compromettere la leggibilit\`a del documento.

\section{Riferimenti}

\subsection{Normativi} % oppure rif. a Norme di progetto con leggi e tutto
\begin{itemize}
  \item ISO/IEC 12207:1995 - Cicli di vita software
  \item ISO/IEC 9126:2001 - Quality Model
  \item \emph{NormeDiProgetto-\versionenormeprogetto.pdf} che regola e
  accompagna tutti i documenti ufficiali.
\end{itemize}
\newpage
\subsection{Informativi}
\begin{itemize}
  \item Capitolato d'appalto CO2-NETMUS del corso di Ingegneria del Software
  A.A. 2010/11 :\\
  \url{http://www.math.unipd.it/~tullio/IS-1/2010/Progetto/NetMus.pdf}
  \item Slide delle lezioni del corso:\\
  \url{http://www.math.unipd.it/~tullio/IS-1/2010/}
  \item Verbale intervista proponente:\\
  \co{allegato Verbale-1.0.pdf}
  \item Sistema di cloud Google App Engine:\\
  \url{http://code.google.com/intl/it/appengine/}
\end{itemize}


% INIZIO CAPITOLO 2

\chapter{Capitolato C1 MindSlide}
\thispagestyle{fancy}
Il capitolato C1 MindSlide (Software di presentazione basato su mind mapping)
proposto dall'azienda Zucchetti richiede l'utilizzo della tecnologia
\underline{HTML5}. Dopo un'attenta lettura della specifica, si ritiene non
sia abbastanza stimolante.

\chapter{Capitolato C2 Netmus}
\thispagestyle{fancy}
Il capitolato C2 NetMus (Libreria musicale personale, virtuale e distribuita)
proposto dal prof Riccardo Cardin richiede l'utilizzo delle tecnologie \underline{Google
App Engine}, \underline{Google Web Toolkit} e \underline{JavaFX}. Dopo
un'accurata lettura e un'intensa discussione si \`e deciso di accettare questo progetto.
\section{Studio del dominio} 
Le conoscenze e le capacit\`a per l'implementazione del prodotto si basano su
questi ambiti:
\begin{itemize}
  \item \underline{internet} e \underline{cloud computing}
  \item gestione remota, concorrenza e distribuzione
  \item programmazione \underline{Java} e utilizzo tecnologie Google App
  Engine, Google Web Toolkit e \underline{RIA}
  \item progettazione con \underline{UML} e \underline{Gantt}
  \item \underline{sistemi operativi} vari \underline{Windows},
  \underline{Linux}, \underline{MacOsx})
\end{itemize}

\section{Fattibilit\`a Tecnologica e Tempistica}
Considerando la preparazione ricevuta ai corsi universitari di Programmazione
Concorrente e Distribuita ed Ingegneria del Software si dovrebbe avere la
capacit\`a per utilizzare i linguaggi Java e UML, e lo
strumento Gantt. Invece per le nuove tecnologie Google
App Engine e Google Web Toolkit si pu\`o utilizzare il corso di
preparazione messo a disposizione dall'azienda Miriade S.p.A. committente del
capitolato C3. Riguardo la piattaforma JavaFx per la realizzazione
delle RIA, non dovrebbe essere un problema recuperare le
informazioni per l'applicazione . Infine dopo aver analizzato la conoscenza interna del gruppo, capacit\`a di apprendimento e impegni vari si dovrebbe riuscire a sviluppare il prodotto e rispettare le date di consegna per le varie revisioni.

\section{Componenti del Progetto e Dettagli sulle Tecnologie}
Il problema che viene posto nel capitolato C2 presenta due principali componenti
da considerare per lo studio della fattibilit\`a. Una componente che gestisce il
recupero delle informazioni dei brani musicali dal dispositivo di riproduzione personale dell'utente.
L'altra componente gestisce la memorizzazzione e visualizzazione della libreria virtuale personale 
dell'utente.

\subsection{Componente di estrazione delle informazioni}
Per sviluppare la prima componente si \`e pensato di ricorrere alla creazione di
una Rich Internet Application, ossia un'\underline{applicazione web} che possiede le
caratteristiche e le funzionalit\`a delle \underline{applicazioni desktop}, senza
per\`o necessitare dell'installazione sul \underline{disco fisso}. Questa scelta
\`e stata dettata dal fatto che le RIA si caratterizzano per la dimensione interattiva, la multimedialit\`a e la velocit\`a d'esecuzione.
La parte dell'applicazione che elabora i dati \`e trasferita a livello client e
fornisce una pronta risposta all'interfaccia utente, mentre la gran parte dei dati e dell'applicazione 
rimane sul server remoto, con notevole alleggerimento per il computer utente. 
Queste caratteristiche rendono l'applicazione il meno invasiva possibile per
l'utente.

\subsection{Componente di memorizzazione e gestione della libreria virtuale}
La seconda componente dovr\`a poter essere accessibile all'utente come
applicazione web. Il sistema consigliato per lo sviluppo, ossia Google App Engine, 
\`e una tecnologia cloud computing che permette l'utilizzo di risorse hardware o
 software distribuite in remoto. User\`a inoltre il database distribuito
 \underline{Google DataStore}. Questo Database utiliza una sintassi di
 linguaggio simile all'\underline{SQL} chiamato \underline{GQL}.

\section{Rischi}
L'impegno nell'analisi della fattibilit\`a dovrebbe aver minimizzato le variabili di rischio.
Tuttavia la mancanza di esperienza o imprevisti inattesi come la mancanza di apporto di risorse umane potrebbero provocare ritardi o addiritttura la non riuscita nel completamento del lavoro.

\chapter{Capitolato C3 RescueMe}
\thispagestyle{fancy}
Il capitolato C3 RescueMe (Sistema globale per gestione catastrofi naturali)
proposto dall'azienda Miriade S.p.A. richiede l'utilizzo delle tecnologie
\underline{Google App Engine}, \underline{Google Web Toolkit} e
\underline{Spring MVC}. Dopo una profonda analisi, si ritiene che gli strumenti
da utilizzare siano di interesse comune al gruppo. Per\`o si pensa che l'implementazione del progetto richieda tempi molto lunghi, col rischio di non rispettare i tempi di consegna prefissati.

\chapter{Capitolato C4 QWAD}
\thispagestyle{fancy}
Il capitolato C4 QWAD (Quick Web Apps Designer) proposto dalla \underline{Onlus}
Informatici Senza Frontiere presuppone la continuazione di un lavoro gi\`a svolto da altri
team di sviluppo. Questo fa ermergere la difficolt\`a di studio e preparazione degli strumenti richiesti. Pur considerando la gratuit\`a sociale di tale progetto, si \`e comunque deciso di non accettarlo.

\end{document}